В рамках выполнения выпускной квалификационной работы был разработан программный комплекс для интерполяции отображений и решений дифференциальных уравнений на основе адаптивных разреженных сеток с поддержкой квадратичного базиса и параллельного построения.

Полный исходный код проекта, включая вычислительный модуль, серверный слой, клиент визуализации и вспомогательные средства развертывания, размещён в открытом репозитории на платформе GitHub.

\textbf{Ссылка на репозиторий:} \url{https://github.com/DesCaldnd/Diplom}

\begin{figure}[H]
    \centering
    \qrcode[height=5cm]{https://github.com/DesCaldnd/Diplom}
    
    \vspace{1ex}
    \caption{QR-код со ссылкой на репозиторий с исходным кодом проекта}
    \label{fig:qr_code_github}
\end{figure}

В состав репозитория входят:

\begin{enumerate}
    \item исходный код вычислительного модуля \texttt{compute}, реализующего адаптивные разреженные сетки;
    \item gRPC-сервис, обеспечивающий удалённый вызов построения сетки и вычисления отображений;
    \item клиентское приложение визуализации на Flutter;
    \item protobuf-описание контракта обмена данными;
    \item набор тестов, примеров и бенчмарков для исследования корректности и производительности алгоритмов.
\end{enumerate}